\section{Proposed}
\label{Pro}

\subsection{Dirty Ancilla AND Gadgets}
\label{sub-synth}
Here we take the dirty ancilla construction in~\ref{sub-mct}, and create a gadget to
realize an AND using dirty ancilla. First, we separate the construction into
two parts, as in~\figref{fig-dirt-mct-org}: we label these \emph{compute}, and \emph{uncompute}.
We observe that the Toffoli gate in compute can be replaced with a single target controlled
gate for an arbitrary function F; that is, it's a gate that flips its target bit when a
function F is true. We can implement this function F as a quantum
subcircuit~\cite{bib-amy-phase-state}, and so we can realize a recursive construction, depicted
in~\figref{fig-compute}. The Toffoli gate in the uncompute gadget can thus be replaced with the
gate F's adjoint. For a relative phase gate, this means implementing F's components' adjoints in
reverse, to realize a gate that implements the same function, but with the opposite relative phase
~\cite{bib-maslov-rtof}.

\begin{figure}[t]
  \centering
  \scalebox{1.0} {
    \begin{quantikz}[row sep=0.5em, column sep=0.5em]
  \setwiretype{n}    &     &                                                                                                                                                                 & \push{Compute} &          &  & \push{Uncompute}                                                                                                                                                     &  &  \\
  \setwiretype{n}    &     &                                                                                                                                                                 &                &          &  &                                                                                                                                                                      &  &  \\  
  \lstick{\ket{x_0}} & \qw & \qw \gategroup[5,steps=3,style={dashed,rounded corners,fill=blue!20, inner xsep=2pt},background,label style={label position=above,anchor=north,yshift=-0.2cm}]{}& \ctrl{1}       & \qw      &  & \ctrl{1} \gategroup[5,steps=2,style={dashed,rounded corners,fill=red!20, inner xsep=2pt},background,label style={label position=above,anchor=north,yshift=-0.2cm}]{} &  & \rstick{\ket{x_0}} \\
  \lstick{\ket{x_1}} & \qw & \qw                                                                                                                                                             & \ctrl{2}       & \qw      &  & \ctrl{2}                                                                                                                                                             &  & \rstick{\ket{x_1}} \\
  \lstick{\ket{x_2}} & \qw & \ctrl{1}                                                                                                                                                        & \qw            & \ctrl{1} &  & \qw                                                                                                                                                                  &  & \rstick{\ket{x_2}} \\
  \lstick{a}         & \qw & \ctrl{1}                                                                                                                                                        & \targ{}        & \ctrl{1} &  & \targ{}                                                                                                                                                              &  & \rstick{a} \\  
  \lstick{\ket{y}}   & \qw & \targ{}                                                                                                                                                         & \qw            & \targ{}  &  & \qw                                                                                                                                                                  &  & \rstick{\ket{x_0 \cdot x_1 \oplus y}}  
\end{quantikz}

%% \begin{quantikz}[row sep=0.5em, column sep=0.5em]
%%   \lstick{\ket{x_0}} & \qw & \qw      & \ctrl{1} & \qw      & \ctrl{1} & \qw & \rstick{\ket{x_0}} \\
%%   \lstick{\ket{x_1}} & \qw & \qw      & \ctrl{2} & \qw      & \ctrl{2} & \qw & \rstick{\ket{x_1}} \\
%%   \lstick{\ket{x_2}} & \qw & \ctrl{1}\gategroup[5,steps=3,style={dashed,rounded corners},label style={above}]{Compute} & \qw      & \ctrl{1} & \qw      & \qw & \rstick{\ket{x_2}} \\
%%   \lstick{a}         & \qw & \ctrl{1} & \targ{}  & \ctrl{1} & \targ{}  &     & \rstick{a} \\  
%%   \lstick{\ket{y}}   & \qw & \targ{}  & \qw      & \targ{}\gategroup[5,steps=1,style={dashed,rounded corners},label style={above}]{Uncompute} & \qw      & \qw & \rstick{\ket{x_0 \cdot x_1 \oplus y}}  \\
%% \end{quantikz}
%% \begin{quantikz}[row sep=0.5em, column sep=0.5em]
%%   \lstick{\ket{x_0}} & \qw & \qw      & \ctrl{1} & \qw      & \ctrl{1} & \qw & \rstick{\ket{x_0}} \\
%%   \lstick{\ket{x_1}} & \qw & \qw      & \ctrl{2} & \qw      & \ctrl{2} & \qw & \rstick{\ket{x_1}} \\
%%   \lstick{\ket{x_2}} & \qw & \ctrl{1} & \qw      & \ctrl{1} & \qw      & \qw & \rstick{\ket{x_2}} \\
%%   \lstick{a}         & \qw & \ctrl{1} & \targ{}  & \ctrl{1} & \targ{}  &  & \rstick{a} \\  
%%   \lstick{\ket{y}}   & \qw & \targ{}  & \qw      & \targ{}  & \qw      & \qw & \rstick{\ket{x_0 \cdot x_1 \oplus y}}  \\
%% \end{quantikz}                        


  }
  \caption{Dirty MCT decomposition base}
  \label{fig-dirt-mct-org}
\end{figure}

\begin{figure}[t]
  \centering
  \scalebox{1.0} {
    \begin{quantikz}[row sep=0.3cm, column sep=0.3cm]
\qw              & \qw             & \qw        & \gate[2]{F}         & \qw               & \qw      & \qw    \\
\qw              & \qw             & \qw        & \ghost{F}\ctrl{2}   & \qw               & \qw      & \qw    \\
\lstick{$G$}     & \qw             & \ctrl{1}   & \qw                 & \ctrl{1}          & \qw      & \qw    \\
\lstick{$a_0$}   & \qw             & \gate{Z}   & \targ{}             & \gate{Z^\dagger}  & \qw      & \qw    \\
\lstick{$t$}     & \gate{H}        & \ctrl{-1}  & \qw                 & \ctrl{-1}         & \gate{H} & \qw    \\
\end{quantikz}

%% \begin{quantikz}[row sep=0.3cm, column sep=0.55cm]
%% \qw              & \qw             & \qw        & \qw             & \qw              & \gate[2]{F_i}       & \qw                    & \qw       & \qw             & \qw      & \qw        & \qw              & \gate[2]{F_i}        & \qw                     & \qw        \\
%% \qw              & \qw             & \qw        & \qw             & \qw              & \ghost{F_i}         & \qw                    & \qw       & \qw             & \qw      & \qw        & \qw              & \ghost{F_i}          & \qw                     & \qw        \\
%% \qw              & \qw             & \qw        & \qw             & \ctrl{2}         & \qw                 & \ctrl{2}               & \qw       & \qw             & \qw      & \qw        & \ctrl{2}         & \qw                 & \ctrl{2}                & \qw         \\
%% \qw              & \qw             & \octrl{2}  & \qw             & \qw              & \qw                 & \qw                    & \qw       & \qw             & \qw      & \qw        & \qw              & \qw                 & \qw                     & \qw         \\
%% \lstick{$a_1$}   & \qw             & \qw        & \qw             & \gate{i\omega Z} & \targ{}             & \gate{i\omega Z^\dagger} & \qw       & \octrl{1}       & \qw      & \qw        & \gate{i\omega Z} & \targ{}             & \gate{i\omega Z^\dagger}  & \qw          \\
%% \lstick{$a_2$}   & \qw             & \gate{iZ}  & \gate{H}        & \ctrl{-1}        & \qw                 & \ctrl{-1}              & \gate{H}  & \gate{iZ^\dagger} & \qw      & \gate{H}   & \ctrl{-1}        & \qw                 & \ctrl{-1}               & \gate{H}     \\
%% \lstick{}        & \gate{H}        & \ctrl{-1}  & \qw             & \qw              & \qw                 & \qw                    & \qw       & \ctrl{-1}       & \gate{H} & \qw        & \qw              & \qw                 & \qw                     & \qw          \\
%%\end{quantikz}


  }
  \caption{Compute gadget}
  \label{fig-compute}
\end{figure}

We utilize the fact that the CCZ has two complementary implementations to our advantage by replacing the left CCZ
gate with~\figref{fig-cnot-t-ccz} and the right with \figref{fig-alt-cnot-t-ccz}. In so doing we can combine their
sum-of-paths to get the below sum-of-paths equation:

\begin{equation}
  \label{eq-toff-tot}
  \begin{aligned}
    &\frac{1}{4}z + \frac{1}{4}a + \frac{1}{4}f - \frac{1}{4}(z \oplus a) \\
    &\qquad -\frac{1}{4}(z \oplus f) - \frac{1}{4}(a \oplus f) + \frac{1}{4}(z \oplus a \oplus f)\\
    &- \frac{1}{4}z - \frac{1}{4}(xy \oplus a) - \frac{1}{4} f + \frac{1}{4}(z \oplus xy \oplus a) \\
    &\qquad + \frac{1}{4}(z \oplus f) + \frac{1}{4}(xy \oplus a \oplus f) - \frac{1}{4}(z \oplus xy \oplus a \oplus f)
  \end{aligned}
\end{equation}

This causes some cancellations, and we are left with:

\begin{equation}
  \label{eq-toff-simp}
  \begin{aligned}
    &\frac{1}{4}a - \frac{1}{4}(z \oplus a) - \frac{1}{4}(xy \oplus a) + \frac{1}{4}(z \oplus xy \oplus a)  \\
    &+ \frac{1}{4}(z \oplus a \oplus f) - \frac{1}{4}(a \oplus f) + \frac{1}{4}(xy \oplus a \oplus f) \\
    &- \frac{1}{4}(z \oplus xy \oplus a \oplus f)
  \end{aligned}
\end{equation}

Observe that the first line in\eqnref{eq-toff-simp} contains only terms that have no dependence on $f$.
That means we can bring them out from the summation to create a relative phase version of the compute
gadget, which we depict in~\figref{fig-dirty-gadget-rp}. This means that we can implement the compute
gadget with the same T-count as a relative phase Toffoli gate. We can similarly formulate an uncompute
gadget that is this gadget's conjugate, in order to uncompute the value for the same cost. This implies
that any intermediate computation using dirty ancilla can be performed with the same cost as a relative
phase Toffoli gate. 

The output node will have to implement all 8 terms in~\eqnref{eq-toff-simp}. This means that the T-count of
implementing the output node gadget is 8, so that the difference in T-count between the clean ancilla and
the dirty ancilla decompositions is dependent on the number of these output nodes, which is the number of
output bits of the Boolean function $m$. This means that for an MCT gate, the difference between clean
ancilla and dirty ancilla is 1 T-gate. We can implement these as two RTOF gates.
\begin{figure}
  \scalebox{0.9} {
    \begin{quantikz}[row sep=0.2em, column sep=0.5em]
\qw                & \gate[2]{A}              & \qw             &                    & \setwiretype{n} & \setwiretype{n} & \setwiretype{n} & \setwiretype{q} & \qw           & \qw          & \qw           & \qw                & \qw           & \gate[2]{A}            & \qw           & \qw                & \qw           & \qw            & \qw           & \qw                     \\
\qw                & \ghost{A}\ctrl{2}        & \qw             &                    & \setwiretype{n} & \setwiretype{n} & \setwiretype{n} & \setwiretype{q} & \qw           & \qw          & \qw           & \qw                & \qw           & \ghost{A}\ctrl{2}      & \qw           & \qw                & \qw           & \qw            & \qw           & \qw                     \\
\setwiretype{n}    & \setwiretype{n}          & \setwiretype{n} &                    & \push{=}        &                 &                 & \setwiretype{n} & \push{\vdots} & \push{\vdots}& \push{\vdots} & \push{\vdots}      & \push{\vdots} & \push{\vdots}          & \push{\vdots} & \push{\vdots}      & \push{\vdots} & \push{\vdots}  & \push{\vdots} & \push{\vdots}               \\
\lstick{$B$}       & \gate[1]{\wedge}\ctrl{2} & \qw             & \push{B}           & \setwiretype{n} & \setwiretype{n} & \setwiretype{n} & \setwiretype{q} & \qw           & \qw          & \targ{}       & \gate{T}           & \ctrl{1}      & \targ[style={double}]{}& \ctrl{1}      & \gate{T^{\dagger}} & \targ{ }      & \qw            & \qw           & \qw   \\
\lstick{$a$}       & \targX{}                 & \qw             & \push{A \oplus a}  & \setwiretype{n} & \setwiretype{n} & \setwiretype{n} & \setwiretype{q} & \qw           & \targ{}      & \ctrl{-1}     & \gate{T^{\dagger}} & \targ{}       & \qw                    & \targ{}       & \gate{T}           & \ctrl{-1}     & \targ{}        & \qw           & \qw   \\
\lstick{$f$}       & \targ[style={double}]{}  & \qw             & \push{AB \oplus f} & \setwiretype{n} & \setwiretype{n} & \setwiretype{n} & \setwiretype{q} & \gate{H}      & \ctrl{-1}    & \qw           & \qw                & \qw           & \qw                    & \qw           & \qw                & \qw           & \ctrl{-1}      & \gate{H}      & \qw         
\end{quantikz}

%% \begin{quantikz}[row sep=0.3cm, column sep=0.55cm]
%% \qw              & \qw             & \qw        & \qw             & \qw              & \gate[2]{F_i}       & \qw                    & \qw       & \qw             & \qw      & \qw        & \qw              & \gate[2]{F_i}        & \qw                     & \qw        \\
%% \qw              & \qw             & \qw        & \qw             & \qw              & \ghost{F_i}         & \qw                    & \qw       & \qw             & \qw      & \qw        & \qw              & \ghost{F_i}          & \qw                     & \qw        \\
%% \qw              & \qw             & \qw        & \qw             & \ctrl{2}         & \qw                 & \ctrl{2}               & \qw       & \qw             & \qw      & \qw        & \ctrl{2}         & \qw                 & \ctrl{2}                & \qw         \\
%% \qw              & \qw             & \octrl{2}  & \qw             & \qw              & \qw                 & \qw                    & \qw       & \qw             & \qw      & \qw        & \qw              & \qw                 & \qw                     & \qw         \\
%% \lstick{$a_1$}   & \qw             & \qw        & \qw             & \gate{i\omega Z} & \targ{}             & \gate{i\omega Z^\dagger} & \qw       & \octrl{1}       & \qw      & \qw        & \gate{i\omega Z} & \targ{}             & \gate{i\omega Z^\dagger}  & \qw          \\
%% \lstick{$a_2$}   & \qw             & \gate{iZ}  & \gate{H}        & \ctrl{-1}        & \qw                 & \ctrl{-1}              & \gate{H}  & \gate{iZ^\dagger} & \qw      & \gate{H}   & \ctrl{-1}        & \qw                 & \ctrl{-1}               & \gate{H}     \\
%% \lstick{}        & \gate{H}        & \ctrl{-1}  & \qw             & \qw              & \qw                 & \qw                    & \qw       & \ctrl{-1}       & \gate{H} & \qw        & \qw              & \qw                 & \qw                     & \qw          \\
%%\end{quantikz}


  }
  \caption{Gadget for one of A or B is a primary input}
  \label{fig-prop-and}
\end{figure}
\begin{figure}
  \scalebox{1.0} {
    %% \begin{quantikz}[row sep=0.25cm, column sep=0.35cm]
%% \lstick{} & \ctrl{2}  & \lstick{} &\ghost{\qw}         & \ghost{\qw} & \qw & \qw                & \targ{}   & \qw      & \ctrl{1} & \gate{T}           & \ctrl{1} & \qw      & \targ{}    & \qw \\
%% \lstick{} & \ctrl{1}  & \lstick{} &\push{=} & \ghost{\qw} & \qw & \qw                & \qw       & \ctrl{1} & \targ{}  & \gate{T^{\dagger}} & \targ{}  & \ctrl{1} & \qw        & \qw \\
%% \lstick{} & \gate{iZ} & \lstick{} &\push{ } & \ghost{\qw} & \qw & \gate{T^{\dagger}} & \ctrl{-2} & \targ{}  & \qw      & \gate{T}           & \qw      & \targ{}  & \ctrl{-2}  & \qw 
%% \end{quantikz}
\begin{quantikz}[row sep=0.25cm, column sep=0.35cm]
\lstick{$\ket{x_a}$} & \ctrl{1}                & \setwiretype{n} &          && \setwiretype{q} & \qw &          & \qw & \qw                & \targ{}   & \qw      & \ctrl{1} & \gate{T}           & \ctrl{1} & \qw      & \targ{}    &          & \qw \\
\lstick{$\ket{x_b}$} & \ctrl{1}                & \setwiretype{n} & \push{=} && \setwiretype{q} & \qw &          & \qw & \qw                & \qw       & \ctrl{1} & \targ{}  & \gate{T^{\dagger}} & \targ{}  & \ctrl{1} & \qw        &          & \qw \\
\lstick{$\ket{x_c}$} & \targ[style={double}]{} & \setwiretype{n} &          && \setwiretype{q} & \qw & \gate{H} & \qw & \gate{T^{\dagger}} & \ctrl{-2} & \targ{}  & \qw      & \gate{T}           & \qw      & \targ{}  & \ctrl{-2}  & \gate{H} & \qw \\
\end{quantikz}


  }\\\hline
  \caption{Gadget for when both A and B are primary inputs}
  \label{fig-prop-rtof}
\end{figure}

\begin{figure}
  \scalebox{0.6} {
    \begin{quantikz}[row sep=0.2em, column sep=0.5em]
\qw                & \gate[2]{A}              & \qw             &                    & \setwiretype{n} & \setwiretype{n} & \setwiretype{n} & \setwiretype{q} & \qw           & \qw          & \qw           & \qw                & \qw           & \gate[2]{A}            & \qw           & \qw                & \qw           & \qw            & \qw           & \qw                     \\
\qw                & \ghost{A}\ctrl{2}        & \qw             &                    & \setwiretype{n} & \setwiretype{n} & \setwiretype{n} & \setwiretype{q} & \qw           & \qw          & \qw           & \qw                & \qw           & \ghost{A}\ctrl{2}      & \qw           & \qw                & \qw           & \qw            & \qw           & \qw                     \\
\setwiretype{n}    & \setwiretype{n}          & \setwiretype{n} &                    & \push{=}        &                 &                 & \setwiretype{n} & \push{\vdots} & \push{\vdots}& \push{\vdots} & \push{\vdots}      & \push{\vdots} & \push{\vdots}          & \push{\vdots} & \push{\vdots}      & \push{\vdots} & \push{\vdots}  & \push{\vdots} & \push{\vdots}               \\
\lstick{$B$}       & \gate[1]{\wedge}\ctrl{2} & \qw             & \push{B}           & \setwiretype{n} & \setwiretype{n} & \setwiretype{n} & \setwiretype{q} & \qw           & \qw          & \targ{}       & \gate{T}           & \ctrl{1}      & \targ[style={double}]{}& \ctrl{1}      & \gate{T^{\dagger}} & \targ{ }      & \qw            & \qw           & \qw   \\
\lstick{$a$}       & \targX{}                 & \qw             & \push{A \oplus a}  & \setwiretype{n} & \setwiretype{n} & \setwiretype{n} & \setwiretype{q} & \qw           & \targ{}      & \ctrl{-1}     & \gate{T^{\dagger}} & \targ{}       & \qw                    & \targ{}       & \gate{T}           & \ctrl{-1}     & \targ{}        & \qw           & \qw   \\
\lstick{$f$}       & \targ[style={double}]{}  & \qw             & \push{AB \oplus f} & \setwiretype{n} & \setwiretype{n} & \setwiretype{n} & \setwiretype{q} & \gate{H}      & \ctrl{-1}    & \qw           & \qw                & \qw           & \qw                    & \qw           & \qw                & \qw           & \ctrl{-1}      & \gate{H}      & \qw         
\end{quantikz}

%% \begin{quantikz}[row sep=0.3cm, column sep=0.55cm]
%% \qw              & \qw             & \qw        & \qw             & \qw              & \gate[2]{F_i}       & \qw                    & \qw       & \qw             & \qw      & \qw        & \qw              & \gate[2]{F_i}        & \qw                     & \qw        \\
%% \qw              & \qw             & \qw        & \qw             & \qw              & \ghost{F_i}         & \qw                    & \qw       & \qw             & \qw      & \qw        & \qw              & \ghost{F_i}          & \qw                     & \qw        \\
%% \qw              & \qw             & \qw        & \qw             & \ctrl{2}         & \qw                 & \ctrl{2}               & \qw       & \qw             & \qw      & \qw        & \ctrl{2}         & \qw                 & \ctrl{2}                & \qw         \\
%% \qw              & \qw             & \octrl{2}  & \qw             & \qw              & \qw                 & \qw                    & \qw       & \qw             & \qw      & \qw        & \qw              & \qw                 & \qw                     & \qw         \\
%% \lstick{$a_1$}   & \qw             & \qw        & \qw             & \gate{i\omega Z} & \targ{}             & \gate{i\omega Z^\dagger} & \qw       & \octrl{1}       & \qw      & \qw        & \gate{i\omega Z} & \targ{}             & \gate{i\omega Z^\dagger}  & \qw          \\
%% \lstick{$a_2$}   & \qw             & \gate{iZ}  & \gate{H}        & \ctrl{-1}        & \qw                 & \ctrl{-1}              & \gate{H}  & \gate{iZ^\dagger} & \qw      & \gate{H}   & \ctrl{-1}        & \qw                 & \ctrl{-1}               & \gate{H}     \\
%% \lstick{}        & \gate{H}        & \ctrl{-1}  & \qw             & \qw              & \qw                 & \qw                    & \qw       & \ctrl{-1}       & \gate{H} & \qw        & \qw              & \qw                 & \qw                     & \qw          \\
%%\end{quantikz}


  }
  \caption{Relative Phase version of dirty ancilla intermediate compute gadget}
  \label{fig-dirty-gadget-rp}
\end{figure}

\subsection{Optimizing }

This allows us to create a a quantum circuit based off the gadgets in~\tabref{fig-gadgets}. This is depicted
in~\figref{fig-mct-gadgets}. 

\begin{figure}[t]
  \centering
  \scalebox{1.0} {
    \begin{quantikz}[row sep=0.2em, column sep=0.5em]
\lstick{$x_1$}                & \ctrl{1}                   & \ctrl{1}                           & \qw               \\
\lstick{$x_2$}                & \ctrl{2}                   & \ctrl{2}                           & \qw               \\
\setwiretype{n}               & \setwiretype{n}            & \setwiretype{n}                    & \setwiretype{n}   \\
\lstick{$x_3$}                & \gate[2]{\wedge}           & \gate[2]{\wedge^{\dagger}}         & \qw                 \\
\lstick{$a_1$}                & \ghost{\wedge}\ctrl{1}     & \ghost{\wedge^{\dagger}}\ctrl{1}   & \qw               \\
\setwiretype{n}               &                            &                                    & \setwiretype{n}   \\
\setwiretype{n}\push{\vdots}  & \push{\vdots}              & \push{\vdots}                      & \setwiretype{n}   \\
\setwiretype{n}               &                            &                                    & \setwiretype{n}   \\
\lstick{$x_n$}                & \gate[2]{\wedge}\ctrl{-1}] & \vqw{-1}\vqw{1}                    & \qw               \\
\lstick{$a_{n-2}$}            & \ghost{\wedge}\ctrl{1}     & \targ[style={double}]{}            & \qw               \\
\lstick{$f$}                  & \targ{}                    & \qw                                &
\end{quantikz}

%% \begin{quantikz}[row sep=0.3cm, column sep=0.55cm]
%% \qw              & \qw             & \qw        & \qw             & \qw              & \gate[2]{F_i}       & \qw                    & \qw       & \qw             & \qw      & \qw        & \qw              & \gate[2]{F_i}        & \qw                     & \qw        \\
%% \qw              & \qw             & \qw        & \qw             & \qw              & \ghost{F_i}         & \qw                    & \qw       & \qw             & \qw      & \qw        & \qw              & \ghost{F_i}          & \qw                     & \qw        \\
%% \qw              & \qw             & \qw        & \qw             & \ctrl{2}         & \qw                 & \ctrl{2}               & \qw       & \qw             & \qw      & \qw        & \ctrl{2}         & \qw                 & \ctrl{2}                & \qw         \\
%% \qw              & \qw             & \octrl{2}  & \qw             & \qw              & \qw                 & \qw                    & \qw       & \qw             & \qw      & \qw        & \qw              & \qw                 & \qw                     & \qw         \\
%% \lstick{$a_1$}   & \qw             & \qw        & \qw             & \gate{i\omega Z} & \targ{}             & \gate{i\omega Z^\dagger} & \qw       & \octrl{1}       & \qw      & \qw        & \gate{i\omega Z} & \targ{}             & \gate{i\omega Z^\dagger}  & \qw          \\
%% \lstick{$a_2$}   & \qw             & \gate{iZ}  & \gate{H}        & \ctrl{-1}        & \qw                 & \ctrl{-1}              & \gate{H}  & \gate{iZ^\dagger} & \qw      & \gate{H}   & \ctrl{-1}        & \qw                 & \ctrl{-1}               & \gate{H}     \\
%% \lstick{}        & \gate{H}        & \ctrl{-1}  & \qw             & \qw              & \qw                 & \qw                    & \qw       & \ctrl{-1}       & \gate{H} & \qw        & \qw              & \qw                 & \qw                     & \qw          \\
%%\end{quantikz}


  }
  \caption{Synthesis of~\figref{fig-mct-xag} using dirty ancilla gadgets}
  \label{fig-mct-gadgets}
\end{figure}

This construction is remarkably similar to the ``V-chain'' that is created when the clean ancilla 
decomposition is used, such as in~\figref{fig-clean-mct}. This demonstrates that the difference in
cost between a dirty ancilla decomposition and a clean ancilla decomposition is the cost of the
final $\wedge$ gadget, which uses one more T-gate than the Toffoli gate that the clean ancilla
decomposition uses.

\subsection{Dirty Ancilla-like MCT Decomposition}
